\documentclass[11pt]{article}
\usepackage[top=2cm]{geometry}

\title{Homework 13.2: Bui Tuong Phong's Biography}
\author{Alejandro Cornejo Pérez}
\date{May 2026}
\setlength{\parskip}{\baselineskip}
\begin{document}

\maketitle
Bui Tuong Phong was born in December 14, 1942 in Hanoi Vietnam. He was a computer scientist that became very important for modern computer graphics with this work. He specialized in the realistic rendering of 3D surfaces. He had a very short life but he made contributions that still remain fundamental in sections like video games, animation, and visual effects.

Phong showed strong academic ability from an early age. He pursued higher education in France, where he studied electrical engineering at the Institut Polytechnique de Grenoble. He did very good academically in France and that's why he earned opportunities to continue his studies in the United States. In the US he studied at the University of Utah which was one of the best places for computer graphics research during the 1970s.

His most influential contribution was the development of the Phong reflection model. This model describes how light interacts with surfaces so that computers can simulate realistic shading. This allowed for the rendering of shiny highlights on objects which was way more realistic than earlier methods. In addition to the reflection model, he also developed Phong shading. This was a technique that interpolates surface normals across polygons to produce smooth shading effects which was a significant improvement over flat shading and Gouraud shading. Today, variations of Phong shading are still used in graphics pipelines for video games and simulations.

Phong’s ideas set the bases for the evolution of rendering algorithms, they are still taught in computer graphics courses worldwide. Even though more advanced models like physically based rendering, the principles introduced by Phong remain relevant. Tragically, Bui Tuong Phong died at the age of 32 from leukemia after completing his doctoral work. He had a short life and career but his contributions influenced decades of research and development in computer graphics.

In conclusion Bui Tuong Phong was a brilliant person that invented innovative models of light reflection and shading. His legacy is still important to shape visual realism seen in modern digital media.
\end{document}
